% BHU PYQ -- Professional Exam Template
\documentclass[11pt,a4paper]{article}

\usepackage[a4paper,top=2.5cm,bottom=2.5cm,left=2.8cm,right=2.8cm,headheight=14pt]{geometry}
\usepackage{amsmath,amssymb}
\usepackage[T1]{fontenc}
\usepackage[utf8]{inputenc}
\usepackage{lmodern}
\usepackage{microtype}
\usepackage{fancyhdr}
\usepackage{enumitem}
\usepackage{listings}
\usepackage{hyperref}
\usepackage{lastpage}
\usepackage{parskip}
\usepackage{tikz}

\hypersetup{colorlinks=true,urlcolor=black,linkcolor=black,
  pdftitle={BVCA-202: Problem Solving AI -- BHU PYQ}}

\pagestyle{fancy}
\fancyhf{}
\renewcommand{\headrulewidth}{0.4pt}
\renewcommand{\footrulewidth}{0.4pt}
\fancyhead[L]{\small   \textit{Banaras Hindu University}}
\fancyhead[R]{\small   \textit{BVCA-202: Problem Solving AI}}
\fancyfoot[C]{\small Page  \thepage\ of \pageref{LastPage}}

\lstset{
  basicstyle=\small tfamily,
  frame=single,
  breaklines=true,
  columns=flexible,
  keepspaces=true
}


\newlist{parts}{enumerate}{2}
\setlist[parts,1]{label=(\alph*),leftmargin=*,itemsep=4pt,topsep=3pt}
\setlist[parts,2]{label=(\roman*),leftmargin=*,itemsep=2pt,topsep=2pt}

\newcommand{\pts}[1]{\hfill   \textit{[#1]}}


\begin{document}

\begin{center}
  {\large\bfseries BANARAS HINDU UNIVERSITY}\\[2pt]
  {
\normalsize B.Voc. Semester II Examination, 2021-22}\\[6pt]
  \rule{0.8\linewidth}{0.5pt}\\[6pt]
  {\large\bfseries Paper: BVCA-202 --- Problem Solving AI}\\[4pt]
  {
\normalsize Time: Three Hours \hfill Maximum Marks: 70}
\end{center}

\rule{\linewidth}{0.4pt}

\smallskip
\smallskip



\noindent   \textit{Instructions: Answer all questions in Part A, any four questions from Part B, and any three questions from Part C.}

\bigskip

\bigskip


\noindent {\large\bfseries PART --- A}
\rule{\linewidth}{0.4pt}\smallskip




\noindent   \textit{Note: Answer all questions. Each question carries 2 marks.} \hfill   \textbf{[10 $ \times$ 2 = 20]}

\medskip


\noindent   \textbf{Question 1.} Answer the following:

\begin{parts}
  \item Briefly explain the Turing Test.
  \item Discuss the Chinese Room Argument.
  \item Describe the PEAS model on which an AI agent works.
  \item Transform the following English sentences into logical notation:
  
\begin{parts}
    \item Not all integers are odd.
    \item Some numbers are not real.
  \end{parts}
  \item Write the following rules of inference related to the language of propositions:
  
\begin{parts}
    \item Modus Ponens
    \item Hypothetical Syllogism
  \end{parts}
  \item Which of the following is not a logic programming language?
  
\begin{parts}
    \item Prolog
    \item ROOP
    \item Janus
    \item C
  \end{parts}
  \item Which of the following is not a heuristic search?
  
\begin{parts}
    \item Best First Search
    \item Hill Climbing
    \item Depth First Search
    \item A* Search
  \end{parts}
  \item Which of the following is not a proposition?
  
\begin{parts}
    \item $X = 3$
    \item Moon is made up of cheese.
    \item $X > 6$
    \item $3 + 3 = 6$
  \end{parts}
  \item Which of the following is false, where the underlying universe of discourse is the set of integers?
  
\begin{parts}
    \item $\forall x [x < x+1]$
    \item $\exists x [x < x+1]$
    \item $\exists x [x = x+1]$
    \item $\forall x [x = x]$
  \end{parts}
  \item Which of the following techniques is not a structured knowledge representation technique?
  
\begin{parts}
    \item Semantic Nets
    \item Production Rules
    \item Scripts
    \item Frames
  \end{parts}
\end{parts}

\bigskip


\noindent {\large\bfseries PART --- B}
\rule{\linewidth}{0.4pt}\smallskip




\noindent   \textit{Note: Answer any four questions from the following. Each question carries 5 marks.} \hfill   \textbf{[4 $ \times$ 5 = 20]}

\medskip


\noindent   \textbf{Question 2.} Give examples of two agents with their PEAS representation.

\medskip


\noindent   \textbf{Question 3.} Explain the Hill Climbing Algorithm along with its types.

\medskip


\noindent   \textbf{Question 4.} If horses fly or cows eat grass, then the mosquito is the national bird. If the mosquito is the national bird, then peanut butter tastes good on hot dogs. But peanut butter tastes terrible on hot dogs. Therefore, cows don't eat grass. Prove or disprove this argument using the resolution principle.
\centerline{   \textbf{OR}}
Prove or disprove the validity of the following argument with respect to the given statements:

\begin{parts}
  \item David is a professor.
  \item All professors are people.
  \item John is the dean.
  \item Deans are professors.
  \item All professors consider the dean a friend or don't know him.
  \item Everyone is a friend of someone.
  \item People only criticize people that are not their friends.
  \item David criticizes John.
  \item   \textbf{Conclusion}: John is not a friend of David.
\end{parts}

\medskip


\noindent   \textbf{Question 5.} Consider the following graph. Find the most cost-effective path to reach from start state $A$ to final state $J$ using the $A^*$ Algorithm:

\begin{center}

\begin{tikzpicture}[>=stealth, every node/.style={draw, circle, minimum size=8mm, inner sep=0pt}]
  
ode (A) at (0,2) {A};
  
ode (B) at (2,4) {B};
  
ode (F) at (2,0) {F};
  
ode (C) at (4,4) {C};
  
ode (G) at (4,0) {G};
  
ode (D) at (6,4.5) {D};
  
ode (H) at (6,0.5) {H};
  
ode (E) at (8,4) {E};
  
ode (I) at (8,0) {I};
  
ode (J) at (10,2) {J};
  \draw[->, thick] (A) -- node[above left] {6} (B);
  \draw[->, thick] (A) -- node[below left] {3} (F);
  \draw[->, thick] (B) -- node[above] {3} (C);
  \draw[->, thick] (B) -- node[above right] {2} (D);
  \draw[->, thick] (C) -- node[above right] {1} (D);
  \draw[->, thick] (C) -- node[above] {5} (E);
  \draw[->, thick] (D) -- node[above right] {8} (E);
  \draw[->, thick] (F) -- node[below] {1} (G);
  \draw[->, thick] (F) -- node[below right] {7} (H);
  \draw[->, thick] (G) -- node[below] {3} (I);
  \draw[->, thick] (H) -- node[below left] {2} (I);
  \draw[->, thick] (I) -- node[below right] {3} (J);
  \draw[->, thick] (I) -- node[right] {5} (E);
\end{tikzpicture}
\end{center}


\begin{center}

\begin{tabular}{|c|c||c|c|}
  \hline
   \textbf{Node} &   \textbf{h(n)} &   \textbf{Node} &   \textbf{h(n)} \\
  \hline
  A & 10 & F & 6 \\
  \hline
  B & 8 & G & 5 \\
  \hline
  C & 5 & H & 3 \\
  \hline
  D & 7 & I & 1 \\
  \hline
  E & 3 & J & 0 \\
  \hline
\end{tabular}
\end{center}
\centerline{   \textbf{OR}}
Write the pseudo-code for the Min-Max algorithm and explain its working with an example.

\medskip


\noindent   \textbf{Question 6.} Write the advantages and disadvantages of the logical representation technique of knowledge representation.
\centerline{   \textbf{OR}}
Write a Prolog program to create the following knowledge base and return the answers for the queries:

\begin{parts}
  \item   \textbf{Knowledge Base}:
  
\begin{itemize}
    \item Lili is happy if she dances.
    \item Tom is hungry if he is searching for food.
    \item Jack and Bill are friends if both of them love to play cricket.
    \item Ryan will go to play if the school is closed and he is free.
  \end{itemize}
  \item   \textbf{Queries}:
  
\begin{itemize}
    \item Is Tom a cat?
    \item Does Kunal love to eat pasta?
    \item Is Lili happy?
    \item Will Ryan go to play?
  \end{itemize}
\end{parts}

\bigskip


\noindent {\large\bfseries PART --- C}
\rule{\linewidth}{0.4pt}\smallskip




\noindent   \textit{Note: Answer any three questions from the following. Each question carries 10 marks.} \hfill   \textbf{[3 $ \times$ 10 = 30]}

\medskip


\noindent   \textbf{Question 7.} Explain different types of Agents defined based on their degree of perceived intelligence and capability.

\medskip


\noindent   \textbf{Question 8.} Explain how the resolution principle is used in the Prolog language by an example.

\medskip


\noindent   \textbf{Question 9.} Write the Best First Search algorithm. What is the difference between Greedy versus $A^*$ Best First Search Algorithm?

\medskip


\noindent   \textbf{Question 10.} What do you understand by structured knowledge representation techniques? Explain any one of these techniques.

\vfill
\rule{\linewidth}{0.4pt}

\begin{center}\small
\href{https://bhu-exam.vercel.app/}{Click here to visit our website}\\[4pt]\href{https://me-aryan.vercel.app/}{Click here to visit our contributor}\\[4pt]\href{https://quashan-me.vercel.app/}{Click here to visit our contributor}
\end{center}
\end{document}
